\documentclass{article} % For LaTeX2e

% Classic Preprint template v1.01; 13.04.2026

% Based on the COLM format, with several design and usability changes.
% Main template and comments updated by Kirill B.
% This template now supports a single public mode: preprint.
% Use the package option [preprint] to show authors.
\usepackage[preprint]{classic_preprint}

% Optional custom header message shown in the running header.
% Uncomment and edit the next line to override the default preprint header.
% Example:
% \paperheadmessage{Preprint. Version 3, April 2026}
% If this command is omitted, the style file uses its built-in default text.
%%\paperheadmessage{Preprint. Version 3, April 2026}

% Accent color used for links and decorative elements.
\usepackage{xcolor}
\definecolor{Maroon}{cmyk}{0, 0.87, 0.68, 0.32}

% Better micro-typography and colored hyperlinks.
\usepackage{microtype}
\usepackage[colorlinks=true, urlcolor=Maroon, linkcolor=Maroon, citecolor=Maroon]{hyperref}
\usepackage{url}

% Table styling and decorative dropped capital support.
\usepackage{booktabs}
\usepackage{lettrine}

% Decorative initial letter font used in the sample text below.
\input GoudyIn.fd
\newcommand*\initfamily{\usefont{U}{GoudyIn}{xl}{n}}

% NOTE ABOUT PAGE GEOMETRY
% The style file already defines the page dimensions.
% Avoid loading geometry in the usual way, because it can silently override
% the margins and text block dimensions chosen by the template.
% If geometry is absolutely required for some advanced purpose, use:
% \usepackage[pass]{geometry}
%
% Current page design from the style file:
% - Paper size: US Letter (8.5in x 11in)
% - Text width: 460pt, which is about 6.36in
% - Text height: 9.0in
% - Left and right margins: about 1.07in each
% - Top margin to the text block: about 1.1in, depending slightly on header spacing
% - Bottom margin: about 0.85in
%
% This produces a page that feels fairly full and substantial, rather than thin,
% while still keeping a classic single-column academic look.

% Optional line-number package.
% This template does not enable line numbers in preprint mode.
% If you personally want them for drafting, uncomment \linenumbers after
% \begin{document}. They are intentionally off by default here.
%\usepackage{lineno}

\hypersetup{colorlinks=true, urlcolor=Maroon, linkcolor=Maroon, citecolor=Maroon}

% Paper title.
% Use a manual line break only if the title is genuinely too long for one line.
\title{Paper Title\\
A subtitle or second line can go here}

% Author block.
% Replace the dummy names, affiliations, and repository URL with your own.
% The style file formats this block automatically in preprint mode.
\author{%
Alice Smith\textsuperscript{$1$} \; Bob Johnson\thanks{Equal contribution. Correspondence to \texttt{alice.smith@example.com}.}\,\,\textsuperscript{$,1,2$} \; Carol Lee\footnotemark[1]\,\,\textsuperscript{$,1$} \; David Brown\footnotemark[1]\,\,\textsuperscript{$,1$} \\\bfseries
Emma Wilson\textsuperscript{$1,2$} \; Frank Miller\textsuperscript{$1$} \; Grace Taylor\textsuperscript{$1$} \; Henry Davis\textsuperscript{$3$} \\\bfseries
Ivy Moore\textsuperscript{$1$} \; Jack Anderson\textsuperscript{$4$} \; Karen Thomas\textsuperscript{$1$} \\
\textsuperscript{$1$}University A \; \textsuperscript{$2$}Institute B \; \textsuperscript{$3$}University C \; \textsuperscript{$4$}Research Lab D
\\[2ex]
\url{https://github.com/example/project}
}

% The \author macro works with any number of authors.
% There are two useful separators for multi-author layouts:
% - \And lets LaTeX decide where to break the line.
% - \AND forces a line break at that position.

% Small margin markers sometimes used while drafting.
\newcommand{\fix}{\marginpar{FIX}}
\newcommand{\new}{\marginpar{NEW}}

% Simple keyword formatting helper.
\newcommand{\kwsep}{\;\textbullet\;}
\newcommand{\keywords}[1]{%
  \par\noindent\textbf{Keywords: }#1
}

\begin{document}

% Uncomment the next line only if you want line numbers for your own draft.
%\linenumbers

% Builds the title, header text, and author block.
\maketitle

\begin{abstract}
This template is called \emph{Classic Preprint}. It is based on the COLM design, but has been simplified into a single preprint-oriented format with cleaner comments and a customizable header message by yours truly, Kirill Bykov. The current layout uses US Letter
paper, a text width of 460pt, and a text height of 9.0in, which gives a full
but still visually balanced single-column page. The side margins are a little
over one inch, which keeps the paper from looking thin, while the overall page
still retains a classic academic appearance.
\end{abstract}

% Optional keyword line below the abstract.
\keywords{paper template\kwsep preprint\kwsep LaTeX\kwsep typography}

\section{About this template}

\lettrine[lines=4]{\initfamily\textcolor{Maroon}{T}}{his} file is a sample
main document for the \emph{Classic Preprint} template. It is intended as a
starting point for your own paper. The style is based on the COLM format, but
this version removes the older mode structure and keeps only a preprint
workflow. In practice, that means the package is loaded with
\verb|\usepackage[preprint]{classic_preprint}| and there is no longer any need
to think about separate submission or final options in this template.

\subsection{Custom header message}

The style file provides a command for changing the short message shown in the
running header:

\begin{center}
\verb|\paperheadmessage{Preprint. Version 3, April 2026}|
\end{center}

Use this command in the preamble, after loading the style package and before
\verb|\begin{document}|. If you do not provide it, the template falls back to
the built-in default header text from the style file. This is useful for labels
such as a version note, technical report note, revision date, or a submission
status line that you want to control directly from the main \texttt{.tex} file.

\subsection{Current page dimensions}

The present design uses US Letter paper, which is \mbox{8.5in \(\times\) 11in}.
The style file sets the text block to 460pt in width, which is about 6.36in,
and 9.0in in height. That leaves side margins of about 1.07in on each side.
Visually, this gives the paper a fairly full page without making the text block
feel excessively wide. The vertical layout is also compact, which helps the
page look substantial and not sparse.

For a classic single-column academic layout, this is a reasonable balance:
there is enough white space to frame the text, but not so much that the page
looks thin or under-filled. If you want a slightly more relaxed look later, the
first things to experiment with would be a slightly narrower text width and a
slightly shorter text height.

\section{How to use the template}

Replace the title, abstract, authors, affiliations, keywords, and body text
with your own material. The decorative drop cap at the beginning of a section
is optional. The table and figure examples below are included only to show how
the page elements look with the current styling.

\subsection{Citations}

This template loads \texttt{natbib} through the style file, so you can use
commands such as \verb|\citet{key}| and \verb|\citep{key}| in the usual way.
Keep the bibliography style consistent across the whole document.

\subsection{Footnotes}

Footnotes appear at the bottom of the page and are styled by the template.
Here is a small example footnote\footnote{This is an example footnote placed in
the sample document.} to demonstrate the appearance.

\subsection{Figures}

Figures should be clear and readable, with captions placed below the figure.
The example below is only a placeholder.

\begin{figure}[t]
\begin{center}
\fbox{\rule[-.5cm]{0cm}{4cm} \rule[-.5cm]{4cm}{0cm}}
\end{center}
\caption{Sample figure caption.}
\end{figure}

\subsection{Tables}

Tables should be neat and easy to read. The example below uses
\texttt{booktabs}, which is usually a good default choice.

\begin{table}[t]
\begin{center}
\begin{tabular}{ll}
\toprule
\multicolumn{1}{c}{\bf PART}  &\multicolumn{1}{c}{\bf DESCRIPTION} \\
\midrule
Dendrite         & Input terminal \\
Axon             & Output terminal \\
Soma             & Cell body (contains cell nucleus) \\
\bottomrule
\end{tabular}
\end{center}
\caption{Sample table title}\label{sample-table}
\end{table}

\section{Practical notes}

Try not to load packages that silently change page geometry, spacing, or font
choices unless you really need them. This template is meant to provide a stable
baseline with a classic look. Most users will only need to edit the title,
author block, header message, abstract, keywords, and paper body.

\section{Summary}

The goal of this sample file is to document the template itself while also
serving as a clean starting point for a new paper. The design keeps the page
visually pleasing, avoids an overly thin text block, and exposes a simple
customization point through \verb|\paperheadmessage{...}|.

\end{document}
